\begin{trapbox}
	\begin{description}[style=nextline, leftmargin=2.8em, labelsep=0.5em, itemsep=0.35em]
		\item[\textbf{\textcolor{CTrap}{易错点一（公式混淆）}}] 将正四面体高误记为底面等边三角形的高$\frac{\sqrt{3}}{2}a$，或记成$\frac{\sqrt{2}}{3}a$等。
		\item[\textbf{\textcolor{CTrap}{正确理解（记清公式）：}}] 正四面体高为$h = \frac{\sqrt{6}}{3}a \approx 0.816a$，底面三角形高约为$0.866a$，两者不同。
		\item[\textbf{\textcolor{CTrap}{错因分析（平面干扰）：}}] 将平面三角形高与空间正四面体高混淆，忽视了正四面体的立体结构。
	\end{description}

	\medskip
	\begin{description}[style=nextline, leftmargin=2.8em, labelsep=0.5em, itemsep=0.35em]
		\item[\textbf{\textcolor{CTrap}{易错点二（适用条件混淆）}}] 把侧棱与底面边长不等的正三棱锥当作正四面体，直接套用高公式。
		\item[\textbf{\textcolor{CTrap}{正确理解（界定图形）：}}] 正四面体要求所有棱长相等，一般正三棱锥（底面正三角形，侧棱与底棱可能不等）的高公式不同。
		\item[\textbf{\textcolor{CTrap}{错因分析（概念混淆）：}}] 概念不清，误以为正三棱锥就是正四面体。
	\end{description}

	\medskip
	\begin{description}[style=nextline, leftmargin=2.8em, labelsep=0.5em, itemsep=0.35em]
		\item[\textbf{\textcolor{CTrap}{易错点三（距离错用）}}] 在推导高时，将底面中心到顶点距离误用为底面中心到边的距离（内切圆半径$\frac{a}{2\sqrt{3}}$），导致结果错误。
		\item[\textbf{\textcolor{CTrap}{正确理解（区分内外径）：}}] 底面中心到顶点距离是外接圆半径$\frac{a}{\sqrt{3}}$，不是内切圆半径。
		\item[\textbf{\textcolor{CTrap}{错因分析（内外混淆）：}}] 混淆了正三角形中心到顶点和中心到边的距离，或者未理解重心分中线2:1的具体应用。
	\end{description}
\end{trapbox}
